\section*{الفصل \ref{ch:b3:complex-kinetics} --- الحركية المعقدة: السلاسل والإنزيمات والتذبذبات}

\begin{solution}{exo:b3:complex-kinetics:1}
$\ce{Cl2 -> 2 Cl}$: بدء؛ $\ce{Cl + H2 -> HCl + H}$ و$\ce{H + Cl2 -> HCl + Cl}$:
انتشار؛ $\ce{2 Cl -> Cl2}$: إنهاء. الحوامل: Cl وH. والتفاعل الإجمالي (مجموع
خطوتي الانتشار): $\ce{H2 + Cl2 -> 2 HCl}$.
\end{solution}

\begin{solution}{exo:b3:complex-kinetics:2}
$[\mathrm M]_{1/2} = k_2/k_{-1} = \qty{1.0e-3}{mol/L} = \qty{1.0}{mol/m^3}$؛ $p = cRT = 1.0 \times 8.314 \times 750 =
\qty{6.2}{kPa}$.
\end{solution}

\begin{solution}{exo:b3:complex-kinetics:3}
$v/V_{\max} = [\mathrm S]/(K_M + [\mathrm S])$: $\frac12$، $\frac34$، $\frac{9}{10}$. وبلوغ 90\,\% من $V_{\max}$ يقتضي $[\mathrm S] = 9K_M$.
\end{solution}

\begin{solution}{exo:b3:complex-kinetics:4}
$k_{\mathrm{cat}}/K_M = 100/\num{0.80e-3} = \qty{1.25e5}{L.mol^{-1}.s^{-1}}$: أي أدنى من حد الانتشار بنحو 60\,000
مرة، وقريب من القيمة النموذجية.
\end{solution}

\begin{solution}{exo:b3:complex-kinetics:5}
الحالة المستقرة من أجل \ce{CH3CO}: $k_2[\ce{CH3}][\ce{CH3CHO}] = k_3[\ce{CH3CO}]$. والحالة المستقرة من أجل \ce{CH3}:
$k_1[\ce{CH3CHO}] - k_2[\ce{CH3}][\ce{CH3CHO}] + k_3[\ce{CH3CO}] - 2k_4[\ce{CH3}]^2 = 0$. وبالجمع: $k_1[\ce{CH3CHO}] =
2k_4[\ce{CH3}]^2$، ومنه $[\ce{CH3}] = (k_1/2k_4)^{1/2}[\ce{CH3CHO}]^{1/2}$، و$\dd[\ce{CH4}]/\dd t = k_2[\ce{CH3}][\ce{CH3CHO}] =
k_2(k_1/2k_4)^{1/2}[\ce{CH3CHO}]^{3/2}$.
\end{solution}

\begin{solution}{exo:b3:complex-kinetics:6}
عند نصف التحويل $[\ce{H2}] = [\ce{Br2}] = \frac12$، $[\ce{HBr}] = 1$ (وحدات نسبية):
$v/v_0 = \frac12 \times (\frac12)^{1/2}/(1 + 0.10 \times 1/0.5) = 0.354/1.2 = 0.29$. ودون تثبيط كانت
ستكون 0.354: فالتثبيط يقسم السرعة على 1.2 إضافية.
\end{solution}

\begin{solution}{exo:b3:complex-kinetics:7}
يحدث الانفجار حين $1000p > 6000/p + 10p^2$، أي $10p^3 - 1000p^2 + 6000 < 0$، وجذراها
الموجبان 2.48 و\qty{99.9}{kPa}: فالخليط ينفجر بين نحو 2.5 
و\qty{100}{kPa} (الحدان الأول والثاني في النموذج).
\end{solution}

\begin{solution}{exo:b3:complex-kinetics:8}
$1/v = 1/V_{\max} + (K_M/V_{\max})/[\mathrm S]$: $5.0 = 1/V_{\max} + 2.0K_M/V_{\max}$ و$2.5 = 1/V_{\max} + 0.5K_M/V_{\max}$.
وبالطرح: $K_M/V_{\max} = \qty{1.667}{s.mM.\micro M^{-1}}$، ثم $1/V_{\max} = 1.667$: $V_{\max} = \qty{0.60}{\micro M.s^{-1}}$،
$K_M = \qty{1.0}{mM}$.
\end{solution}

\begin{solution}{exo:b3:complex-kinetics:9}
المقطع نفسه على محور $1/v$: تنافسي. $K_M^{\mathrm{app}} = 1/0.417 = \qty{2.40}{mM} = 3K_M$، ومنه
$1 + 50/K_i = 3$ و$K_i = \qty{25}{\micro M}$.
\end{solution}

\begin{solution}{exo:b3:complex-kinetics:10}
$N = \qty{1.01}{mol/L}$: $t_{1/2} = \ln(1.01/0.010 - 1)/(0.50 \times 1.01) = \ln100/0.505 = \qty{9.1}{s}$. 
والسرعة $kx(N - x)$ أعظمية عند $x = N/2$: أي في اللحظة نفسها، وهي نقطة انعطاف المنحنى
ذي شكل S.
\end{solution}

\begin{solution}{exo:b3:complex-kinetics:11}
الأثر $B - 2$، والمحدد 1. $B = 1.5$: $\lambda = -0.25 \pm 0.968\iu$، بؤرة مستقرة (تذبذبات
متخامدة نحو الحالة المستقرة). $B = 2.5$: $\lambda = 0.25 \pm 0.968\iu$، بؤرة غير مستقرة:
فتنمو التذبذبات حتى تبلغ \omterm{def:b3:complex-kinetics:oscillating}{الدورة الحدية}.
\end{solution}

\begin{solution}{exo:b3:complex-kinetics:12}
$1/\tau = \num{1.0e8} \times \num{5.40e-5} + \num{1.0e3} = \qty{6.4e3}{s^{-1}}$، $\tau = \qty{156}{\micro s}$. قِس $\tau$ عند عدة
تراكيز كلية: فالمقدار $1/\tau$ بدلالة $[\mathrm A]_e + [\mathrm B]_e$ مستقيم ميله $k_1$ 
ومقطعه $k_{-1}$.
\end{solution}

\begin{solution}{pb:b3:complex-kinetics:1}
\textbf{1.} في البداية يكون $[\mathrm S]$ معلومًا، والناتج، الذي قد يثبّط أو يتفاعل
عكسيًا، غائبًا.
\textbf{2.} نقاط $[\mathrm S]$ المنخفضة، عند قيم $1/[\mathrm S]$ و$1/v$ الكبيرة: فهي تحدد الميل، 
والقلب يضخّم أخطاءها.
\textbf{3.} لأنه يرجّح كل سرعة مقيسة كما قيست؛ أما مستقيم لاينويفر--بيرك فيعطي
وسائط منحازة.
\textbf{4.} $K_M$: القيمة $0.773$ أدنى من 0.801 بمقدار 1.2 خطأ معياري؛ $V_{\max}$: القيمة $0.491$ أدنى من 0.502 بمقدار 2.2
خطأ معياري. وكلتاهما منخفضة، كما يُتوقع من الانحياز.
\textbf{5.} $k_{\mathrm{cat}} = 0.502/0.0050 = \qty{100}{s^{-1}}$.
\textbf{6.} $100/\num{0.801e-3} = \qty{1.25e5}{L.mol^{-1}.s^{-1}}$.
\textbf{7.} أدنى من حد الانتشار بنحو 60\,000 مرة؛ فهو \omterm{def:b3:complex-kinetics:enzyme}{إنزيم} نموذجي.
\textbf{8.} تنافسي: $V_{\max}$ لا تتغير، و$K_M$ الظاهري يكبر مع $[\mathrm I]$.
\textbf{9.} $K_M^{\mathrm{app}} = K_M(1 + [\mathrm I]/K_i) = K_M + (K_M/K_i)[\mathrm I]$.
\textbf{10.} $K_i = 0.799/0.0321 = \qty{24.9}{\micro M}$.
\textbf{11.} $24.9 \pm 4.30 \times 0.9 = 24.9 \pm \qty{3.9}{\micro M}$.
\textbf{12.} ثابت تفكك معقد الإنزيم--المثبط: أي تركيز
المثبط الذي يشغل نصف \omterm{def:b3:complex-kinetics:enzyme}{الإنزيم} الحر.
\textbf{13.} $v_i/v_0 = (K_M + [\mathrm S])/(K_M(1 + [\mathrm I]/K_i) + [\mathrm S])$.
\textbf{14.} $2/3 = 0.67$.
\textbf{15.} $2/12 = 0.17$.
\textbf{16.} $11/21 = 0.52$: عند تركيز مرتفع للركيزة تتغلب الركيزة في التنافس على
المثبط.
\textbf{17.} عند $[\mathrm S] = K_M$، تعطي $v_i/v_0 = 2/(2 + [\mathrm I]/K_i) = 0.10$ العلاقة $[\mathrm I]/K_i = 18$.
\textbf{18.} $18 \times 24.9 = \qty{448}{\micro M}$.
\textbf{19.} الخطوة الأولى سريعة: فكل جزيء \omterm{def:b3:complex-kinetics:enzyme}{إنزيم} يحرر بسرعة جزيئًا واحدًا من $\mathrm P_1$
ويُحتجز في صورة الأسيل--إنزيم؛ وبعد ذلك لا يظهر $\mathrm P_1$ إلا بالسرعة التي تحرر بها
الخطوة الثانية البطيئة \omterm{def:b3:complex-kinetics:enzyme}{الإنزيم}.
\textbf{20.} $1.28/2.0 = 0.64 = (k_2/(k_2 + k_3))^2$، ومنه $k_2/(k_2 + k_3) = 0.80$.
\textbf{21.} $200/2.0 = \qty{100}{s^{-1}}$، أي القيمة نفسها كما في الجزء الأول.
\textbf{22.} $k_2 = 0.80 \times 625 = \qty{500}{s^{-1}}$، $k_3 = \qty{125}{s^{-1}}$.
\textbf{23.} $500 \times 125/625 = \qty{100}{s^{-1}}$.
\textbf{24.} \textbf{$[\mathrm I]_{90\,\%} = 18K_i \approx \qty{0.45}{mM}$ عند $[\mathrm S] = K_M$.}
\end{solution}
